\documentclass[11pt,fontset=fandol]{ctexart}
\usepackage[a4paper,margin=23mm]{geometry}
\usepackage{amsmath,amssymb,mathtools,bm,mathrsfs,esint}
\usepackage{graphicx,longtable,booktabs,array,calc}
\usepackage[normalem]{ulem}
\usepackage{xcolor,hyperref}
\hypersetup{hidelinks}
\setcounter{secnumdepth}{-1}
\newcounter{none}
\providecommand{\tightlist}{\setlength{\itemsep}{0pt}\setlength{\parskip}{0pt}}
\providecommand{\passthrough}[1]{#1}
\usepackage{amssymb,mathrsfs,esint,bm,newunicodechar}
\allowdisplaybreaks
\setlength{\emergencystretch}{3em}
\providecommand{\llangle}{\langle\!\langle}
\providecommand{\rrangle}{\rangle\!\rangle}
\newunicodechar{∞}{\ensuremath{\infty}}
\newunicodechar{∈}{\ensuremath{\in}}
\newunicodechar{∑}{\ensuremath{\sum}}
\newunicodechar{∫}{\ensuremath{\int}}
\newunicodechar{≤}{\ensuremath{\leq}}
\newunicodechar{≥}{\ensuremath{\geq}}
\newunicodechar{≠}{\ensuremath{\neq}}
\newunicodechar{→}{\ensuremath{\to}}
\newunicodechar{⇒}{\ensuremath{\Rightarrow}}
\newunicodechar{α}{\ensuremath{\alpha}}
\newunicodechar{β}{\ensuremath{\beta}}
\newunicodechar{γ}{\ensuremath{\gamma}}
\newunicodechar{λ}{\ensuremath{\lambda}}
\newunicodechar{φ}{\ensuremath{\varphi}}
\newunicodechar{ε}{\ensuremath{\varepsilon}}
\newunicodechar{σ}{\ensuremath{\sigma}}
\newunicodechar{θ}{\ensuremath{\theta}}
\newunicodechar{δ}{\ensuremath{\delta}}
\newunicodechar{η}{\ensuremath{\eta}}
\newunicodechar{π}{\ensuremath{\pi}}
\newunicodechar{ξ}{\ensuremath{\xi}}
\newunicodechar{Σ}{\ensuremath{\Sigma}}
\newunicodechar{Δ}{\ensuremath{\Delta}}
\newunicodechar{∂}{\ensuremath{\partial}}
\newunicodechar{∀}{\ensuremath{\forall}}
\newunicodechar{∃}{\ensuremath{\exists}}
\newunicodechar{√}{\ensuremath{\surd}}
\newunicodechar{−}{\ensuremath{-}}
\newunicodechar{μ}{\ensuremath{\mu}}
\newunicodechar{ρ}{\ensuremath{\rho}}
\newunicodechar{τ}{\ensuremath{\tau}}
\newunicodechar{ω}{\ensuremath{\omega}}
\newunicodechar{Ω}{\ensuremath{\Omega}}
\newunicodechar{χ}{\ensuremath{\chi}}
\newunicodechar{ψ}{\ensuremath{\psi}}
\newunicodechar{Γ}{\ensuremath{\Gamma}}
\newunicodechar{Φ}{\ensuremath{\Phi}}
\newunicodechar{Λ}{\ensuremath{\Lambda}}
\newunicodechar{ℝ}{\ensuremath{\mathbb{R}}}
\newunicodechar{ℤ}{\ensuremath{\mathbb{Z}}}
\newunicodechar{ℕ}{\ensuremath{\mathbb{N}}}
\newunicodechar{△}{\ensuremath{\triangle}}
\newunicodechar{⊗}{\ensuremath{\otimes}}
\begin{document}
\section{2019年杭州电子科技大学601数学分析考研真题参考答案}\label{ux5e74ux676dux5ddeux7535ux5b50ux79d1ux6280ux5927ux5b66601ux6570ux5b66ux5206ux6790ux8003ux7814ux771fux9898ux53c2ux8003ux7b54ux6848}

来源：\texttt{HDU/601数学分析/2019/2019年杭州电子科技大学601数学分析考研真题参考答案.pdf}。原
PDF 页序：1, 2, 3, 4, 5, 6, 7, 8, 9。

\section{2019年数学分析打印参考答案：OCR核定整理}\label{ux5e74ux6570ux5b66ux5206ux6790ux6253ux5370ux53c2ux8003ux7b54ux6848ocrux6838ux5b9aux6574ux7406}

\begin{quote}
标注``【OCR公式整理】''的受损段，由助手按同题题面和可核算推导恢复，不作为原作者逐字转录；原卷符号存在歧义时保留待核对说明。
\end{quote}

\section{杭州电子科技大学2019年数学分析真题解析}\label{ux676dux5ddeux7535ux5b50ux79d1ux6280ux5927ux5b662019ux5e74ux6570ux5b66ux5206ux6790ux771fux9898ux89e3ux6790}

\subsection{一.填空题（本大蔬英6小题，每小题5分，本大题共40分）}\label{ux4e00.ux586bux7a7aux9898ux672cux5927ux852cux82f16ux5c0fux9898ux6bcfux5c0fux98985ux5206ux672cux5927ux9898ux517140ux5206}

1.极限
\(\lim _ { n \to \infty } ( { \frac { 1 } { n + { \sqrt { 1 } } } } + { \frac { 1 } { n + { \sqrt { 2 } } } } + \cdots + { \frac { 1 } { n + { \sqrt { n } } } } ) = (\)

【宝哥解答】

\[
\frac { 1 } { n + \sqrt { 1 } } + \frac { 1 } { n + \sqrt { 2 } } + \cdots + \frac { 1 } { n + \sqrt { n } } \leq \frac { n } { n + \sqrt { 1 } }  1 , n  \infty ,
\]

\[
{ \frac { 1 } { n + { \sqrt { 1 } } } } + { \frac { 1 } { n + { \sqrt { 2 } } } } + \cdots + { \frac { 1 } { n + { \sqrt { n } } } } \geq { \frac { n } { n + { \sqrt { n } } } } \to 1 , n \to \infty ,
\]

故將
\(\lim _ { n \to \infty } ( { \frac { 1 } { n + { \sqrt { 1 } } } } + { \frac { 1 } { n + { \sqrt { 2 } } } } + \cdots + { \frac { 1 } { n + { \sqrt { n } } } } ) = 1\)

\begin{enumerate}
\def\labelenumi{\arabic{enumi}.}
\setcounter{enumi}{1}
\tightlist
\item
  设
  \(g(0)=g'(0)=0\)，\(f(x)=g(x)\sin(1/x)\)（\(x\ne0\)），\(f(0)=0\)，求
  \(f'(0)\)。
\end{enumerate}

【OCR公式整理】有 \(g(x)/x\to g'(0)=0\)，且 \(|\sin(1/x)|\le1\)，所以 \[
f'(0)=\lim_{x\to0}\frac{g(x)}x\sin(1/x)=0.
\]

3.曲线
\(L : \left\{ { \begin{array} { l } { x = 1 - t ^ { 2 } } \\ { y = t - t ^ { 2 } } \end{array} } \right.\)
在t=1所对应的法线方程为（

\subsection{【宝哥解答】}\label{ux5b9dux54e5ux89e3ux7b54}

\(t \neq 0\) 时
\(\frac { d y } { d x } = \frac { \displaystyle \frac { d y } { d t } } { \displaystyle \frac { d x } { d t } } = \frac { 1 - 2 t } { - 2 t } = 1 - \frac { 1 } { 2 t } \ , \frac { d y } { d x } \big | _ { t = 1 } = 1 - \frac { 1 } { 2 } = \frac { 1 } { 2 } . t = 1\)
时 ， \(x = 0\)

\(L : \left\{ \begin{array} { l l } { x = 1 - t ^ { 2 } } \\ { y = t - t ^ { 2 } } \end{array} \right.\)
在t=1所对应的法线方程为 \(y = - 2 x\) \(y = 0\) 故曲线

\begin{enumerate}
\def\labelenumi{\arabic{enumi}.}
\setcounter{enumi}{3}
\tightlist
\item
  曲线 \(y=x^4-6x^3+12x^2-x+1\) 的拐点。
\end{enumerate}

【OCR公式整理】\(y''=12(x-1)(x-2)\)，在 \(1,2\) 两侧均改变符号。故拐点为
\((1,7)\)、\((2,15)\)。

\begin{enumerate}
\def\labelenumi{\arabic{enumi}.}
\setcounter{enumi}{4}
\tightlist
\item
  连续函数满足 \(2f(x)+x=1+\int_0^1f(x)\,dx\)，求该积分。
\end{enumerate}

【OCR公式整理】记 \(J=\int_0^1f(x)\,dx\)，两端在 \([0,1]\) 积分，得
\(2J+1/2=1+J\)，所以 \(J=1/2\)。

\begin{enumerate}
\def\labelenumi{\arabic{enumi}.}
\setcounter{enumi}{5}
\tightlist
\item
  设 \(F(y)=\int_y^{y^2}\sin(xy)/x\,dx\)，求 \(F'(y)\)。
\end{enumerate}

【OCR公式整理】换元 \(u=xy\)，有 \[
F(y)=\int_{y^2}^{y^3}\frac{\sin u}{u}\,du,\qquad
F'(y)=\frac{3\sin(y^3)-2\sin(y^2)}y\quad(y\ne0).
\] 若将被积函数在 \(x=0\) 连续延拓，则 \(F'(0)=0\)。

\begin{enumerate}
\def\labelenumi{\arabic{enumi}.}
\setcounter{enumi}{6}
\tightlist
\item
  交换 \(\int_1^{e^2}dy\int_{\ln y}^2f(x,y)\,dx\) 的积分次序。
\end{enumerate}

【OCR公式整理】区域为 \(0\le x\le2,\ 1\le y\le e^x\)，所以结果为 \[
\int_0^2dx\int_1^{e^x}f(x,y)\,dy.
\]

8.若【OCR公式整理】设 \(R=|a|\)；通常半径参数取
\(a>0\)。圆周关于坐标置换及原点反射对称，长度为 \(2\pi R\)，从而 \[
\int_Lx^2\,ds=\int_Ly^2\,ds=\int_Lz^2\,ds
=\frac13R^2(2\pi R)=\frac{2\pi R^3}{3},
\quad\int_Lx\,ds=\int_Ly\,ds=0.
\] 因此原积分为 \(2\pi R^3+2\pi R\)，在 \(a>0\) 的通常约定下为
\(2\pi a^3+2\pi a\)。

\subsection{二.计算题（本大题共5小题，每小题6分，本大题共30分)}\label{ux4e8c.ux8ba1ux7b97ux9898ux672cux5927ux9898ux51715ux5c0fux9898ux6bcfux5c0fux98986ux5206ux672cux5927ux9898ux517130ux5206}

1.求极限值
\(\lim _ { x  0 ^ { + } } ( \frac { \sin x } { x } ) ^ { \frac { 1 } { x ^ { 2 } } }\)

【OCR公式整理】由
\(\ln(\sin x/x)=-x^2/6+o(x^2)\)，取对数再取指数，所求极限为
\(e^{-1/6}\)。

\[
y = - y e ^ { x } + 2 e ^ { y } \sin x + 2 x
\]

\[
y ^ { \prime } ( 0 )
\]

【宝哥解答】

令x=0，可得y=0.

令 \(F ( x , y ) = y + y e ^ { x } - 2 e ^ { y } \sin x - 2 x\) ，则

\[
F _ { x } ( x , y ) = y e ^ { x } - 2 e ^ { y } \cos x - 2 , F _ { y } ( x , y ) = 1 + e ^ { x } - 2 e ^ { y } \sin x
\]

由隐函数存在定理的结论，有

\[
y ^ { \prime } ( 0 ) = - \frac { F _ { x } ( 0 , 0 ) } { F _ { y } ( 0 , 0 ) } = - \frac { - 4 } { 2 } = 2
\]

\begin{enumerate}
\def\labelenumi{\arabic{enumi}.}
\setcounter{enumi}{2}
\tightlist
\item
  设 \(f(t)=te^{-t^2}\)（\(t\ge0\)），\(f(t)=0\)（\(t<0\)），求
  \(\int_1^x f(u-1)\,du\)。
\end{enumerate}

【OCR公式整理】令 \(t=u-1\)，分别按积分上限的位置计算： \[
\int_1^x f(u-1)\,du=
\begin{cases}
0,&x<1,\\
\frac12(1-e^{-(x-1)^2}),&x\ge1.
\end{cases}
\] 原识别在后一分支错误带入的负号由直接求导核定。

4.求级数
\(\sum _ { n = 1 } ^ { \infty } n { \frac { 2 ^ { \frac { n } { 2 } } } { 3 ^ { n } } }\)
的和

【OCR公式整理】对 \(|t|<1\)，几何级数逐项求导给出 \[
\sum_{n=1}^\infty nt^n
=t\left(\sum_{n=0}^\infty t^n\right)'
=\frac{t}{(1-t)^2}.
\] 取 \(t=\sqrt2/3\)，原数项级数之和为 \[
\frac{\sqrt2/3}{(1-\sqrt2/3)^2}
=\frac{3\sqrt2}{(3-\sqrt2)^2}.
\]

5.计算 \(\iint _ { D } ( x + y + 1 ) ^ { 2 } d x d y\) ，其中
\(D = \left\{ ( x , y ) \left. 1 \leq x ^ { 2 } + y ^ { 2 } \leq 4 \right. \right\}\)

\subsection{【宝哥解答】}\label{ux5b9dux54e5ux89e3ux7b54-1}

【OCR公式整理】环域 \(D\) 关于两坐标轴对称，故 \(x,y,xy\)
项积分均为零。利用极坐标， \[
\iint_D(x+y+1)^2\,dx\,dy
=2\pi\int_1^2(r^2+1)r\,dr
=\frac{15\pi}{2}+3\pi=\frac{21\pi}{2}.
\]

\subsection{三.综合题（本大题共3题，每题10分，本大题共30分）}\label{ux4e09.ux7efcux5408ux9898ux672cux5927ux9898ux51713ux9898ux6bcfux989810ux5206ux672cux5927ux9898ux517130ux5206}

\begin{enumerate}
\def\labelenumi{\arabic{enumi}.}
\tightlist
\item
  求 \(a\)，使 \[
  f(x,y)=\begin{cases}
  (x^2+y^2)^a\cos\frac1{|x|+|y|},&(x,y)\ne(0,0),\\
  0,&(x,y)=(0,0)
  \end{cases}
  \] 在原点可微。
\end{enumerate}

【OCR公式整理】必要性：沿正 \(x\) 轴，偏导差商为
\(x^{2a-1}\cos(1/x)\)。取 \(x_n=1/(n\pi)\)，若
\(a\le1/2\)，该差商不能有有限极限。故必须 \(a>1/2\)。
充分性：此时两偏导均为零，而且 \[
\frac{|f(x,y)|}{\sqrt{x^2+y^2}}
\le(x^2+y^2)^{a-1/2}\to0.
\] 所以所求条件恰为 \(a>1/2\)。

\begin{enumerate}
\def\labelenumi{\arabic{enumi}.}
\setcounter{enumi}{1}
\tightlist
\item
  讨论 \(\sum_{n=1}^\infty\frac{(x+1)^n}{(n+1)2^n}\)
  的绝对收敛、条件收敛与发散范围。
\end{enumerate}

【OCR公式整理】以 \(x+1\) 为变量的收敛半径为 \(2\)，故在 \(-3<x<1\)
绝对收敛。\(x=-3\) 时为交错级数，条件收敛；\(x=1\)
时为调和型级数，发散。区间外通项不趋于零。因此绝对收敛区间是
\((-3,1)\)，条件收敛点是 \(-3\)，其余点发散。

3.计算曲线积分
\(\int _ { L } { \frac { ( e ^ { x } - 2 x y ) d x - ( x ^ { 2 } + y ^ { 2 } ) d y } { ( x - 1 ) ^ { 2 } + y ^ { 2 } } }\)
其中L是在圆周 \(y = \sqrt { 2 x - x ^ { 2 } }\) 上由
点O(0,0)到点A(1,1)的一段弧. 【OCR公式整理】圆弧上
\((x-1)^2+y^2=1\)，故分母恒为 \(1\)。原微分形式为 \[
(e^x-2xy)\,dx-(x^2+y^2)\,dy
=d\left(e^x-x^2y-\frac{y^3}{3}\right).
\] 从 \(O(0,0)\) 到 \(A(1,1)\) 积分，应同时代入两个端点： \[
\left[e^x-x^2y-\frac{y^3}{3}\right]_{(0,0)}^{(1,1)}
=e-\frac43-1=e-\frac73.
\]

\subsection{四.证明题(本大题共5题，每题10分，本大题共50分)}\label{ux56db.ux8bc1ux660eux9898ux672cux5927ux9898ux51715ux9898ux6bcfux989810ux5206ux672cux5927ux9898ux517150ux5206}

1.设【OCR公式整理】由单调递减， \[
a_{n+1}-a_n=f(n+1)-\int_n^{n+1}f(x)\,dx\le0.
\] 又 \[
a_n=f(n)+\sum_{i=1}^{n-1}\left[f(i)-\int_i^{i+1}f(x)\,dx\right]\ge f(n)>0.
\] 所以数列递减且有下界，因而收敛。

\begin{enumerate}
\def\labelenumi{\arabic{enumi}.}
\setcounter{enumi}{1}
\tightlist
\item
  设 \(f\) 在 \([a,b]\) 二阶可导，图像与端点连线还相交于
  \(C(c,f(c))\)，\(a<c<b\)。第（1）问要求两侧各有一点导数相等。
\end{enumerate}

【OCR整理及待核对】原题第（2）问在真题和答案OCR中均写为``存在
\(f'(\xi)=0\)''，但按这一字面结论命题不成立：取 \(f(x)=x\)
即为反例。二阶可导条件及原解答对 \(f'\) 使用 Rolle
定理的过程支持目标应为
\(f''(\xi)=0\)；原卷该阶数仍需核对。以下保留并证明可以核算的两个结论。

记端点连线斜率为 \(m=(f(b)-f(a))/(b-a)\)。共线条件给出 \[
\frac{f(c)-f(a)}{c-a}=\frac{f(b)-f(c)}{b-c}=m.
\] 分别对 \([a,c]\)、\([c,b]\) 使用 Lagrange 中值定理，得 \[
\xi_1\in(a,c),\quad\xi_2\in(c,b),\qquad f'(\xi_1)=f'(\xi_2)=m.
\] 再对 \(f'\) 在 \([\xi_1,\xi_2]\) 使用 Rolle 定理，得
\(\xi\in(\xi_1,\xi_2)\) 满足 \(f''(\xi)=0\)。这不能推出原OCR所写的
\(f'(\xi)=0\)。

\begin{enumerate}
\def\labelenumi{\arabic{enumi}.}
\setcounter{enumi}{2}
\tightlist
\item
  设 \(1<a<b\)，证明 \(F(x)=\int_e^\infty\cos t/t^x\,dt\) 在 \([a,b]\)
  连续且可微。
\end{enumerate}

【OCR公式整理】对于 \(x\in[a,b]\)， \[
\left|\frac{\cos t}{t^x}\right|\le t^{-a},\qquad
\left|\frac{\partial}{\partial x}\frac{\cos t}{t^x}\right|
\le\frac{\ln t}{t^a}.
\] 两个控制函数在 \([e,\infty)\)
可积。因此原积分及导数积分均一致收敛，故 \(F\) 连续，且 \[
F'(x)=-\int_e^\infty\frac{\cos t\ln t}{t^x}\,dt.
\]

4.设 \(\sum _ { n = 1 } ^ { \infty } a _ { n }\) 绝对收敛，证明
\(\sum _ { n = 1 } ^ { \infty } a _ { n } ( a _ { 1 } + a _ { 2 } + \cdots + a _ { n } )\)
也绝对收敛.

\subsection{【宝哥解答】}\label{ux5b9dux54e5ux89e3ux7b54-2}

\(\sum _ { n = 1 } ^ { \infty } a _ { n }\) 绝对收敛，故
\(\sum _ { n = 1 } ^ { \infty } a _ { n }\)
收敛，故其部分和有界，即存在正常数M，使得

对任意正整数n，都有
\(\left| a _ { 1 } + a _ { 2 } + \cdots + a _ { n } \right| \leq M\)
，故

\[
\left| a _ { n } ( a _ { 1 } + a _ { 2 } + \cdots + a _ { n } ) \right| \leq M | a _ { n } { \big | } , n = 1 , 2 , \cdots
\]

\(\sum _ { n = 1 } ^ { \infty } a _ { n }\) 绝对收敛，故
\(\sum _ { n = 1 } ^ { \infty } \left| a _ { n } \right|\)
收敛.由比较判别法，
\(\sum _ { n = 1 } ^ { \infty } \left| a _ { n } { \big ( } a _ { 1 } + a _ { 2 } + \cdots + a _ { n } { \big ) } \right|\)
收敛.即
\(\sum _ { n = 1 } ^ { \infty } a _ { n } ( a _ { 1 } + a _ { 2 } + \cdots + a _ { n } )\)
也绝对收敛.

5.设 \(R > 0 ; \Omega\) 为曲面 \(z = R ^ { 2 } - x ^ { 2 } - y ^ { 2 }\)
与平面z=0所围成的空间闭区域，Σ为Ω表面的外侧，V为Ω的体积.证明

\[
\iint x ^ { 2 } y e ^ { \sin z } d y d z - x y ^ { 2 } e ^ { \sin z } d z d x + z ( 1 + 2 x y z ) d x d y = V
\]

\subsection{【宝哥解答】}\label{ux5b9dux54e5ux89e3ux7b54-3}

【OCR公式整理】按题面，向量场为 \[
\mathbf F=(x^2ye^{\sin z},-xy^2e^{\sin z},z(1+2xyz)).
\] 其散度为 \[
2xye^{\sin z}-2xye^{\sin z}+1+4xyz=1+4xyz.
\] 立体 \(\Omega=\{0\le z\le R^2-x^2-y^2\}\) 关于 \(x=0\)、\(y=0\)
对称，故 \(\iiint_\Omega xyz\,dV=0\)。由 Gauss 公式，外向通量等于 \[
\iiint_\Omega(1+4xyz)\,dV=\iiint_\Omega1\,dV=V.
\]
\end{document}
